% --- Archon physics formalization source begin ---
% archon:physics
% archon:covers PhyXMiniProblems/problem_phyx_mini_0458.lean
% archon:source-report reports/phyx_mini/problem_phyx_mini_0458.source.json

\chapter{Physics problem phyx\_mini\_0458}
\label{ch:PhyXMiniProblems_problem_phyx_mini_0458}

\paragraph{Problem source.}
\#\# Physical scenario

An insulated cylinder is divided into two parts of \$1 \textbackslash{}, \textbackslash{}text\{m\}\textasciicircum{}3\$ each by an initially locked piston, as shown in figure. Side \$A\$ has air at \$200 \textbackslash{}, \textbackslash{}text\{kPa\}\$, \$300 \textbackslash{}, \textbackslash{}text\{K\}\$, and side \$B\$ has air at \$1.0 \textbackslash{}, \textbackslash{}text\{MPa\}\$, \$1000 \textbackslash{}, \textbackslash{}text\{K\}\$. The piston is now unlocked so that it is free to move, and it conducts heat so that the air comes to a uniform temperature \$T\_A = T\_B\$.

\#\# Question

Find the final \$P\$.

\#\# Answer choices

- A: A: 400 kPa

- B: B: 1000 kPa

- C: C: 500 kPa

- D: D: 613 kPa

\#\# Figure caption (auxiliary; use the image as primary evidence)

The image depicts a cross-sectional view of a container divided into two compartments labeled "A" and "B," both filled with air. A movable, vertical piston separates these compartments. 

- **Compartments**: Both sides are filled with air and labeled accordingly.
- **Piston**: Located in between the compartments, vertically aligned, and depicted with lines indicating it can move up or down.
- **Text**: The compartments are marked with "A" or "B" followed by the word "Air" to indicate the content.
- **Background and Border**: The container is outlined and surrounded by a hatched pattern.

The setup suggests the context of a physics or engineering problem, likely related to pressure or volume changes.

\#\# Dataset metadata

Ideal Gases and Kinetic Theory; Physical Model Grounding Reasoning, Multi-Formula Reasoning

\paragraph{Recorded answer/context.}
D: D: 613 kPa

\paragraph{Figure/image path.}
/root/proposal\_for\_physic/science-mango/phyx\_mini\_run/phyx\_data/test\_image/458.png

\paragraph{Formalization target.}
create a compiling Lean file with sorry bodies at `PhyXMiniProblems/problem\_phyx\_mini\_0458.lean`. The Lean declarations must preserve the physical quantities, dimensions or dimensional roles, figure labels, governing-law hypotheses, and final relation expressed by this problem.
Use Mathlib/Physlib names found through LeanExplore where available. If a physics API is missing, introduce faithful local abstractions rather than scalar placeholder aliases.

\begin{theorem}[Physics formalization target]
\label{thm:physics:phyx_mini_0458:target}
\lean{PhyXMiniProblems.ProblemPhyXMini0458.finalPressureAndEnergyCharacterization}
\uses{def:physics:phyx-mini-0458:phyxminiproblems-problemphyxmini0458-conductingpistoncylindersetup, def:physics:phyx-mini-0458:phyxminiproblems-problemphyxmini0458-finalpressureinkilopascals, def:physics:phyx-mini-0458:phyxminiproblems-problemphyxmini0458-finaltemperatureinkelvins, def:physics:phyx-mini-0458:phyxminiproblems-problemphyxmini0458-statespecificinternalenergyinkilojoulesperkilogram, def:physics:phyx-mini-0458:phyxminiproblems-problemphyxmini0458-matchessuppliedcylinderfigure, def:physics:phyx-mini-0458:phyxminiproblems-problemphyxmini0458-matchesproblemdata, def:physics:phyx-mini-0458:phyxminiproblems-problemphyxmini0458-hasphysicalparameters, def:physics:phyx-mini-0458:phyxminiproblems-problemphyxmini0458-satisfiesconductingpistonidealairlaws, def:physics:phyx-mini-0458:phyxminiproblems-problemphyxmini0458-displayedpressureinkilopascals}
For the insulated two-compartment cylinder, the common final pressure is
\(5T_f/6\) kilopascals, and conservation of internal energy gives the stated
mass-weighted caloric relation.  Consequently the displayed \(613\,\mathrm
{kPa}\) is equivalent to \(T_f=3678/5\,\mathrm K\).
\end{theorem}
\begin{proof}
The initial ideal-gas equations give
\(m_AR=200/300=2/3\) and \(m_BR=1000/1000=1\), so the trapped mass ratio is
\(m_A:m_B=2:3\).  At the final common pressure and temperature, summing the
two ideal-gas equations and using total volume \(2\,\mathrm{m^3}\) gives
\[
 2P_f=(m_A+m_B)RT_f=\frac53T_f,
\]
hence \(P_f=5T_f/6\).  Thermal equilibrium makes both final specific internal
energies equal; multiplying total-energy conservation by the common mass
scale gives
\[
 5u_f=2u_{A,i}+3u_{B,i}.
\]
Finally \(5T_f/6=613\) is equivalent, by nonzero rational cancellation, to
\(T_f=3678/5\).
\end{proof}

% --- Archon declaration topology begin ---
\section{Declaration topology}
\begin{definition}[Volume Quantity]
  \label{def:physics:phyx-mini-0458:phyxminiproblems-problemphyxmini0458-volumequantity}
  \lean{PhyXMiniProblems.ProblemPhyXMini0458.VolumeQuantity}
  Physical volume, carrying dimension L³
\end{definition}
\begin{proof}
  This is the declared model definition of Volume Quantity.
\end{proof}
\begin{definition}[Air Mass Quantity]
  \label{def:physics:phyx-mini-0458:phyxminiproblems-problemphyxmini0458-airmassquantity}
  \lean{PhyXMiniProblems.ProblemPhyXMini0458.AirMassQuantity}
  Physical mass of one trapped air sample
\end{definition}
\begin{proof}
  This is the declared model definition of Air Mass Quantity.
\end{proof}
\begin{definition}[Specific Internal Energy Quantity]
  \label{def:physics:phyx-mini-0458:phyxminiproblems-problemphyxmini0458-specificinternalenergyquantity}
  \lean{PhyXMiniProblems.ProblemPhyXMini0458.SpecificInternalEnergyQuantity}
  Mass-specific internal energy, carrying dimension L² T⁻²
\end{definition}
\begin{proof}
  This is the declared model definition of Specific Internal Energy Quantity.
\end{proof}
\begin{definition}[Specific Gas Constant Quantity]
  \label{def:physics:phyx-mini-0458:phyxminiproblems-problemphyxmini0458-specificgasconstantquantity}
  \lean{PhyXMiniProblems.ProblemPhyXMini0458.SpecificGasConstantQuantity}
  The mass-specific gas constant, with dimension L² T⁻² Θ⁻¹, or energy per unit mass per absolute temperature
\end{definition}
\begin{proof}
  This is the declared model definition of Specific Gas Constant Quantity.
\end{proof}
\begin{definition}[volume In Cubic Meters]
  \label{def:physics:phyx-mini-0458:phyxminiproblems-problemphyxmini0458-volumeincubicmeters}
  \lean{PhyXMiniProblems.ProblemPhyXMini0458.volumeInCubicMeters}
  \uses{def:physics:phyx-mini-0458:phyxminiproblems-problemphyxmini0458-volumequantity}
  Read a physical volume in coherent SI cubic metres
\end{definition}
\begin{proof}
  This is the declared model definition of volume In Cubic Meters.
\end{proof}
\begin{definition}[air Mass In Kilograms]
  \label{def:physics:phyx-mini-0458:phyxminiproblems-problemphyxmini0458-airmassinkilograms}
  \lean{PhyXMiniProblems.ProblemPhyXMini0458.airMassInKilograms}
  \uses{def:physics:phyx-mini-0458:phyxminiproblems-problemphyxmini0458-airmassquantity}
  Read a physical mass in coherent SI kilograms
\end{definition}
\begin{proof}
  This is the declared model definition of air Mass In Kilograms.
\end{proof}
\begin{definition}[pressure In Pascals]
  \label{def:physics:phyx-mini-0458:phyxminiproblems-problemphyxmini0458-pressureinpascals}
  \lean{PhyXMiniProblems.ProblemPhyXMini0458.pressureInPascals}
  Read a physical pressure in pascals
\end{definition}
\begin{proof}
  This is the declared model definition of pressure In Pascals.
\end{proof}
\begin{definition}[pressure In Kilopascals]
  \label{def:physics:phyx-mini-0458:phyxminiproblems-problemphyxmini0458-pressureinkilopascals}
  \lean{PhyXMiniProblems.ProblemPhyXMini0458.pressureInKilopascals}
  \uses{def:physics:phyx-mini-0458:phyxminiproblems-problemphyxmini0458-pressureinpascals}
  Read a physical pressure in kilopascals
\end{definition}
\begin{proof}
  This is the declared model definition of pressure In Kilopascals.
\end{proof}
\begin{definition}[temperature In Kelvins]
  \label{def:physics:phyx-mini-0458:phyxminiproblems-problemphyxmini0458-temperatureinkelvins}
  \lean{PhyXMiniProblems.ProblemPhyXMini0458.temperatureInKelvins}
  Read an absolute temperature in kelvins. Physlib temperatures store a value relative to an explicit zero-preserving temperature unit
\end{definition}
\begin{proof}
  This is the declared model definition of temperature In Kelvins.
\end{proof}
\begin{definition}[specific Internal Energy In Kilojoules Per Kilogram]
  \label{def:physics:phyx-mini-0458:phyxminiproblems-problemphyxmini0458-specificinternalenergyinkilojoulesperkilogram}
  \lean{PhyXMiniProblems.ProblemPhyXMini0458.specificInternalEnergyInKilojoulesPerKilogram}
  \uses{def:physics:phyx-mini-0458:phyxminiproblems-problemphyxmini0458-specificinternalenergyquantity}
  Read specific internal energy in kilojoules per kilogram
\end{definition}
\begin{proof}
  This is the declared model definition of specific Internal Energy In Kilojoules Per Kilogram.
\end{proof}
\begin{definition}[specific Gas Constant In Kilojoules Per Kilogram Kelvin]
  \label{def:physics:phyx-mini-0458:phyxminiproblems-problemphyxmini0458-specificgasconstantinkilojoulesperkilogramkelvin}
  \lean{PhyXMiniProblems.ProblemPhyXMini0458.specificGasConstantInKilojoulesPerKilogramKelvin}
  \uses{def:physics:phyx-mini-0458:phyxminiproblems-problemphyxmini0458-specificgasconstantquantity}
  Read the specific gas constant in kilojoules per kilogram-kelvin
\end{definition}
\begin{proof}
  This is the declared model definition of specific Gas Constant In Kilojoules Per Kilogram Kelvin.
\end{proof}
\begin{definition}[Compartment]
  \label{def:physics:phyx-mini-0458:phyxminiproblems-problemphyxmini0458-compartment}
  \lean{PhyXMiniProblems.ProblemPhyXMini0458.Compartment}
  The compartment labels visible in the supplied figure
\end{definition}
\begin{proof}
  This is the declared model definition of Compartment.
\end{proof}
\begin{definition}[Process Stage]
  \label{def:physics:phyx-mini-0458:phyxminiproblems-problemphyxmini0458-processstage}
  \lean{PhyXMiniProblems.ProblemPhyXMini0458.ProcessStage}
  The stages before and after releasing the piston
\end{definition}
\begin{proof}
  This is the declared model definition of Process Stage.
\end{proof}
\begin{definition}[Gas Species]
  \label{def:physics:phyx-mini-0458:phyxminiproblems-problemphyxmini0458-gasspecies}
  \lean{PhyXMiniProblems.ProblemPhyXMini0458.GasSpecies}
  The gas named in each compartment
\end{definition}
\begin{proof}
  This is the declared model definition of Gas Species.
\end{proof}
\begin{definition}[Piston Constraint]
  \label{def:physics:phyx-mini-0458:phyxminiproblems-problemphyxmini0458-pistonconstraint}
  \lean{PhyXMiniProblems.ProblemPhyXMini0458.PistonConstraint}
  Whether piston translation along the cylinder is constrained
\end{definition}
\begin{proof}
  This is the declared model definition of Piston Constraint.
\end{proof}
\begin{definition}[Figure Object]
  \label{def:physics:phyx-mini-0458:phyxminiproblems-problemphyxmini0458-figureobject}
  \lean{PhyXMiniProblems.ProblemPhyXMini0458.FigureObject}
  Individually identifiable objects in the supplied raster
\end{definition}
\begin{proof}
  This is the declared model definition of Figure Object.
\end{proof}
\begin{definition}[Supplied Two Compartment Cylinder Figure]
  \label{def:physics:phyx-mini-0458:phyxminiproblems-problemphyxmini0458-suppliedtwocompartmentcylinderfigure}
  \lean{PhyXMiniProblems.ProblemPhyXMini0458.SuppliedTwoCompartmentCylinderFigure}
  \uses{def:physics:phyx-mini-0458:phyxminiproblems-problemphyxmini0458-compartment, def:physics:phyx-mini-0458:phyxminiproblems-problemphyxmini0458-figureobject}
  Qualitative information read from the primary image. These fields encode only visible topology, orientation, labels, and colour; no numerical answer occurs in the raster
\end{definition}
\begin{proof}
  This is the declared model definition of Supplied Two Compartment Cylinder Figure.
\end{proof}
\begin{definition}[Compartment Air State]
  \label{def:physics:phyx-mini-0458:phyxminiproblems-problemphyxmini0458-compartmentairstate}
  \lean{PhyXMiniProblems.ProblemPhyXMini0458.CompartmentAirState}
  \uses{def:physics:phyx-mini-0458:phyxminiproblems-problemphyxmini0458-volumequantity, def:physics:phyx-mini-0458:phyxminiproblems-problemphyxmini0458-airmassquantity}
  One equilibrium state of the trapped gas on one side of the piston
\end{definition}
\begin{proof}
  This is the declared model definition of Compartment Air State.
\end{proof}
\begin{definition}[Ideal Air Caloric Model]
  \label{def:physics:phyx-mini-0458:phyxminiproblems-problemphyxmini0458-idealaircaloricmodel}
  \lean{PhyXMiniProblems.ProblemPhyXMini0458.IdealAirCaloricModel}
  \uses{def:physics:phyx-mini-0458:phyxminiproblems-problemphyxmini0458-specificinternalenergyquantity}
  An abstract caloric model for ideal air. The source does not supply its table entries or a closed formula, so the temperature-to-internal-energy relation is kept abstract rather than calibrated to the recorded answer
\end{definition}
\begin{proof}
  This is the declared model definition of Ideal Air Caloric Model.
\end{proof}
\begin{definition}[Conducting Piston Cylinder Setup]
  \label{def:physics:phyx-mini-0458:phyxminiproblems-problemphyxmini0458-conductingpistoncylindersetup}
  \lean{PhyXMiniProblems.ProblemPhyXMini0458.ConductingPistonCylinderSetup}
  \uses{def:physics:phyx-mini-0458:phyxminiproblems-problemphyxmini0458-specificgasconstantquantity, def:physics:phyx-mini-0458:phyxminiproblems-problemphyxmini0458-compartment, def:physics:phyx-mini-0458:phyxminiproblems-problemphyxmini0458-processstage, def:physics:phyx-mini-0458:phyxminiproblems-problemphyxmini0458-gasspecies, def:physics:phyx-mini-0458:phyxminiproblems-problemphyxmini0458-pistonconstraint, def:physics:phyx-mini-0458:phyxminiproblems-problemphyxmini0458-suppliedtwocompartmentcylinderfigure, def:physics:phyx-mini-0458:phyxminiproblems-problemphyxmini0458-compartmentairstate, def:physics:phyx-mini-0458:phyxminiproblems-problemphyxmini0458-idealaircaloricmodel}
  The two-compartment experiment and its independently stored observables
\end{definition}
\begin{proof}
  This is the declared model definition of Conducting Piston Cylinder Setup.
\end{proof}
\begin{definition}[final Pressure]
  \label{def:physics:phyx-mini-0458:phyxminiproblems-problemphyxmini0458-finalpressure}
  \lean{PhyXMiniProblems.ProblemPhyXMini0458.finalPressure}
  \uses{def:physics:phyx-mini-0458:phyxminiproblems-problemphyxmini0458-conductingpistoncylindersetup}
  The common final pressure, named using the left compartment readout
\end{definition}
\begin{proof}
  This is the declared model definition of final Pressure.
\end{proof}
\begin{definition}[final Pressure In Kilopascals]
  \label{def:physics:phyx-mini-0458:phyxminiproblems-problemphyxmini0458-finalpressureinkilopascals}
  \lean{PhyXMiniProblems.ProblemPhyXMini0458.finalPressureInKilopascals}
  \uses{def:physics:phyx-mini-0458:phyxminiproblems-problemphyxmini0458-pressureinkilopascals, def:physics:phyx-mini-0458:phyxminiproblems-problemphyxmini0458-conductingpistoncylindersetup, def:physics:phyx-mini-0458:phyxminiproblems-problemphyxmini0458-finalpressure}
  The requested final-pressure readout in kilopascals
\end{definition}
\begin{proof}
  This is the declared model definition of final Pressure In Kilopascals.
\end{proof}
\begin{definition}[final Temperature In Kelvins]
  \label{def:physics:phyx-mini-0458:phyxminiproblems-problemphyxmini0458-finaltemperatureinkelvins}
  \lean{PhyXMiniProblems.ProblemPhyXMini0458.finalTemperatureInKelvins}
  \uses{def:physics:phyx-mini-0458:phyxminiproblems-problemphyxmini0458-temperatureinkelvins, def:physics:phyx-mini-0458:phyxminiproblems-problemphyxmini0458-conductingpistoncylindersetup}
  The final-temperature readout in kelvins, named using compartment A
\end{definition}
\begin{proof}
  This is the declared model definition of final Temperature In Kelvins.
\end{proof}
\begin{definition}[state Specific Internal Energy In Kilojoules Per Kilogram]
  \label{def:physics:phyx-mini-0458:phyxminiproblems-problemphyxmini0458-statespecificinternalenergyinkilojoulesperkilogram}
  \lean{PhyXMiniProblems.ProblemPhyXMini0458.stateSpecificInternalEnergyInKilojoulesPerKilogram}
  \uses{def:physics:phyx-mini-0458:phyxminiproblems-problemphyxmini0458-specificinternalenergyinkilojoulesperkilogram, def:physics:phyx-mini-0458:phyxminiproblems-problemphyxmini0458-compartment, def:physics:phyx-mini-0458:phyxminiproblems-problemphyxmini0458-processstage, def:physics:phyx-mini-0458:phyxminiproblems-problemphyxmini0458-conductingpistoncylindersetup}
  The ideal-air table readout at one state's temperature
\end{definition}
\begin{proof}
  This is the declared model definition of state Specific Internal Energy In Kilojoules Per Kilogram.
\end{proof}
\begin{definition}[Matches Supplied Cylinder Figure]
  \label{def:physics:phyx-mini-0458:phyxminiproblems-problemphyxmini0458-matchessuppliedcylinderfigure}
  \lean{PhyXMiniProblems.ProblemPhyXMini0458.MatchesSuppliedCylinderFigure}
  \uses{def:physics:phyx-mini-0458:phyxminiproblems-problemphyxmini0458-conductingpistoncylindersetup}
  Exact qualitative topology and labels visible in the supplied raster
\end{definition}
\begin{proof}
  This is the declared model definition of Matches Supplied Cylinder Figure.
\end{proof}
\begin{definition}[Matches Problem Data]
  \label{def:physics:phyx-mini-0458:phyxminiproblems-problemphyxmini0458-matchesproblemdata}
  \lean{PhyXMiniProblems.ProblemPhyXMini0458.MatchesProblemData}
  \uses{def:physics:phyx-mini-0458:phyxminiproblems-problemphyxmini0458-volumeincubicmeters, def:physics:phyx-mini-0458:phyxminiproblems-problemphyxmini0458-pressureinkilopascals, def:physics:phyx-mini-0458:phyxminiproblems-problemphyxmini0458-temperatureinkelvins, def:physics:phyx-mini-0458:phyxminiproblems-problemphyxmini0458-conductingpistoncylindersetup}
  Numerical readouts and apparatus facts stated in the problem. No final pressure, final temperature, air-table entry, or answer label is a field
\end{definition}
\begin{proof}
  This is the declared model definition of Matches Problem Data.
\end{proof}
\begin{definition}[Has Physical Parameters]
  \label{def:physics:phyx-mini-0458:phyxminiproblems-problemphyxmini0458-hasphysicalparameters}
  \lean{PhyXMiniProblems.ProblemPhyXMini0458.HasPhysicalParameters}
  \uses{def:physics:phyx-mini-0458:phyxminiproblems-problemphyxmini0458-volumeincubicmeters, def:physics:phyx-mini-0458:phyxminiproblems-problemphyxmini0458-airmassinkilograms, def:physics:phyx-mini-0458:phyxminiproblems-problemphyxmini0458-pressureinkilopascals, def:physics:phyx-mini-0458:phyxminiproblems-problemphyxmini0458-temperatureinkelvins, def:physics:phyx-mini-0458:phyxminiproblems-problemphyxmini0458-specificgasconstantinkilojoulesperkilogramkelvin, def:physics:phyx-mini-0458:phyxminiproblems-problemphyxmini0458-conductingpistoncylindersetup}
  Positivity conditions selecting physically meaningful gas states
\end{definition}
\begin{proof}
  This is the declared model definition of Has Physical Parameters.
\end{proof}
\begin{definition}[Satisfies Conducting Piston Ideal Air Laws]
  \label{def:physics:phyx-mini-0458:phyxminiproblems-problemphyxmini0458-satisfiesconductingpistonidealairlaws}
  \lean{PhyXMiniProblems.ProblemPhyXMini0458.SatisfiesConductingPistonIdealAirLaws}
  \uses{def:physics:phyx-mini-0458:phyxminiproblems-problemphyxmini0458-volumeincubicmeters, def:physics:phyx-mini-0458:phyxminiproblems-problemphyxmini0458-airmassinkilograms, def:physics:phyx-mini-0458:phyxminiproblems-problemphyxmini0458-pressureinkilopascals, def:physics:phyx-mini-0458:phyxminiproblems-problemphyxmini0458-temperatureinkelvins, def:physics:phyx-mini-0458:phyxminiproblems-problemphyxmini0458-specificgasconstantinkilojoulesperkilogramkelvin, def:physics:phyx-mini-0458:phyxminiproblems-problemphyxmini0458-conductingpistoncylindersetup, def:physics:phyx-mini-0458:phyxminiproblems-problemphyxmini0458-statespecificinternalenergyinkilojoulesperkilogram}
  Macroscopic governing physics for the two trapped air samples: every state obeys the mass-specific ideal-gas equation PV = mRT; the piston preserves the mass of each trapped sample; the cylinder preserves total volume; insulation preserves the total internal energy of the two air samples; a free conducting piston gives final mechanical and thermal equilibrium. These laws contain no solved pressure, final temperature, air-table calibration, or answer-choice value
\end{definition}
\begin{proof}
  This is the declared model definition of Satisfies Conducting Piston Ideal Air Laws.
\end{proof}
\begin{definition}[Answer Choice]
  \label{def:physics:phyx-mini-0458:phyxminiproblems-problemphyxmini0458-answerchoice}
  \lean{PhyXMiniProblems.ProblemPhyXMini0458.AnswerChoice}
  The four answer labels printed with the source problem
\end{definition}
\begin{proof}
  This is the declared model definition of Answer Choice.
\end{proof}
\begin{definition}[displayed Pressure In Kilopascals]
  \label{def:physics:phyx-mini-0458:phyxminiproblems-problemphyxmini0458-displayedpressureinkilopascals}
  \lean{PhyXMiniProblems.ProblemPhyXMini0458.displayedPressureInKilopascals}
  \uses{def:physics:phyx-mini-0458:phyxminiproblems-problemphyxmini0458-answerchoice}
  The pressure displayed by each choice, in kilopascals
\end{definition}
\begin{proof}
  This is the declared model definition of displayed Pressure In Kilopascals.
\end{proof}
\begin{definition}[recorded Dataset Answer]
  \label{def:physics:phyx-mini-0458:phyxminiproblems-problemphyxmini0458-recordeddatasetanswer}
  \lean{PhyXMiniProblems.ProblemPhyXMini0458.recordedDatasetAnswer}
  \uses{def:physics:phyx-mini-0458:phyxminiproblems-problemphyxmini0458-answerchoice}
  Dataset metadata recording the supplied answer label; never a premise
\end{definition}
\begin{proof}
  This is the declared model definition of recorded Dataset Answer.
\end{proof}
% --- Archon declaration topology end ---
% --- Archon physics formalization source end ---
